\documentclass[11pt,a4paper]{article}
\usepackage[T1]{fontenc}
\usepackage[utf8]{inputenc}
\usepackage{lmodern}
\usepackage[margin=2.2cm]{geometry}
\usepackage{microtype}
\usepackage{enumitem}
\usepackage{xcolor}
\usepackage{hyperref}
\hypersetup{colorlinks=true,linkcolor=blue!50!black,urlcolor=blue!50!black,
  pdftitle={LM-IDNet: Bare Minimum Detector Software}}
\setlength{\parindent}{0pt}
\setlength{\parskip}{0.55em}
\setlist{leftmargin=1.6em,itemsep=0.2em,topsep=0.2em}
\setlength{\emergencystretch}{2em}
\Urlmuskip=0mu plus 2mu

\title{\textbf{LM-IDNet: Bare Minimum Detector Software}\\[0.4em]
\large The remaining software work, in implementation order}
\author{}
\date{15 October 2026}

\begin{document}
\maketitle

\section{What the delivered detector must do}

\textbf{Deliver one working adaptive IoT anomaly detector.} It learns a device's
normal traffic using a Dirichlet--multinomial model and the LM likelihood
implementation, scores subsequent traffic, emits understandable anomaly events,
updates its normal model conservatively, and runs on one Linux edge gateway.
These are the software promises retained from the original LM-IDNet brief and
the proposal's guarded-adaptation design.

Use the existing TCP, ordinary UDP, SSDP, and ARP counts and the configured
ten-minute windows. Train and calibrate on a development machine if convenient;
copy the resulting model and threshold to the gateway. A local collector can
supply completed count windows. The existing PCAP preprocessing path provides
the file-input route. The smallest deployment supports supplied windows; a live
packet collector is necessary only if the delivered interface explicitly
promises to collect packets itself.

\textbf{Reference:} user-confirmed \texttt{master}, commit
\texttt{96fb241433d1d30bd63ddcc16cb38bd27282d1d6}. The original requirements below
describe the remaining software changes against that snapshot and retain the
detector-only scope. The detailed gateway breakdown added on 15 October 2026
uses the current \texttt{feature/edge} implementation. Earlier requirements are
retained as requirements, rather than declarations that every check has passed.

\section{Reuse what already works}

Keep the existing package, configuration files, native LM library, JSON
artefacts, and Docker setup. The following functionality is already present:
\begin{itemize}
  \item Device-specific PCAP preprocessing and validated count windows.
  \item Initial model fitting and calibration of an anomaly threshold.
  \item \texttt{score\_window()}, which produces the score, anomaly decision,
  severity, observed counts, expected proportions, and category residuals.
  \item Periodic alpha refitting with bounded buffers, convergence checks,
  model fingerprints, and update records.
\end{itemize}
Those components provide the detector's core. The four tasks below complete
their validation, update safety, saved state, and deployable execution.

\section{The implementation order}

\subsection{1. Make the existing detector reject incompatible inputs}

\textbf{In plain language:} make sure the detector uses the right device's counts
with the model trained for those counts.

Record the trained window duration and packet-counting policy with the model,
and include them in compatibility checks. The reference model does not store
its window duration. Reuse the existing device, category-order, score-type, and
model--threshold checks; add the missing device check when calibration and
scoring load processed inputs. Reject invalid counts, incompatible durations,
missing required artefacts, and non-finite threshold values. Keep genuine silent
windows distinct from missing observations. Recreate the model and matching
threshold once if their stored format changes.

\textbf{Where:} model/threshold schemas, \path{src/lm_idnet/artifacts.py},
calibration, scoring, and their existing tests.

\textbf{Done when:} valid inputs still produce the existing decisions, and a
wrong device, wrong window duration, invalid threshold, or mismatched
model--threshold pair fails with a clear error.

\subsection{2. Make the existing adaptive updates conservative}

\textbf{In plain language:} let the detector learn ordinary changes without
letting an obvious attack change what it considers normal.

Keep the existing periodic-refit mode. Score each window before considering it
for learning. Keep one bounded buffer of high-confidence accepted windows;
exclude flagged and near-threshold windows. Add one simple upper packet-count
limit derived from benign calibration traffic to keep obvious traffic floods
out of this buffer. This uses the existing total count.

Calculate candidate thresholds from accepted windows, including any fallback
path. The reference implementation uses all recent scores for threshold updates;
a rejected refit must leave the active model and threshold together unchanged.
Retain positive/finite parameter checks and convergence checks. Use fixed,
documented limits on changes to alpha concentration, category proportions, and
the threshold; reject excessive changes. Recalculate the matching threshold
under each candidate model before activating it. Record acceptance or the reason
for rejection.

\textbf{Where:} \path{src/lm_idnet/models/scoring_stage.py}, adaptation
configuration, and focused scoring tests. Reuse the current refit algorithm.

\textbf{Done when:} a valid candidate can update alpha and its threshold;
flagged or excessive-volume inputs cannot enter the learning buffer; and an
invalid, non-converged, or excessive candidate leaves the detector running on
its previous valid state. Conservative gating reduces contamination risk; it
does not establish immunity to every attack.

\subsection{3. Make a restart preserve the accepted model}

\textbf{In plain language:} restarting the program should recover the accepted
detector, and an unwanted update should be reversible.

Save each accepted model and its matching threshold as one validated version.
Keep the previous valid version and a small record identifying the active pair.
Write and validate the new files before atomically switching that record. If a
write or validation fails, retain the previous active pair. On restart, load the
accepted pair; let an operator restore the previous pair. Empty learning buffers
can be rebuilt after restart, so persisting the entire traffic history is
unnecessary. Keep update records and stored versions bounded.

\textbf{Where:} existing artefact helpers and schemas, adaptive scoring, and
the CLI's loading/rollback path.

\textbf{Done when:} restart restores the accepted pair, rollback restores the
previous pair, and an interrupted save cannot mix a new model with an old
threshold. Verify those state transitions with small file-based tests.

\subsection{4. Provide one bounded runtime on the edge gateway}

\textbf{In plain language:} provide one command that keeps reading traffic
windows and immediately writes the detector's decisions.

Add a thin runtime entry point inside the existing package. Load the accepted
model, threshold, and LM backend once. Read one validated count window at a time,
call the existing scoring function, apply the guarded update policy, and emit
one JSON record per scored window. JSON lines on standard input and output are
sufficient. The runtime must accept new observations without requiring a
predeclared list of capture dates. Keep only the active state and bounded
learning buffers in memory; the current batch scorer retains every result.

Reuse the existing event fields as the human-readable evidence. Device and time
identify the alert, while severity and category residuals explain the deviation.
Preserve raw scores in the output. The existing basic metric command can remain
a utility for supplied labels: accuracy and recall already exist; add
\(\mathrm{FPR}=\mathrm{FP}/(\mathrm{FP}+\mathrm{TN})\), returning an unavailable
value if there are no benign labels. Recall is TPR. Labels are never required
for runtime decisions or updates.

Package the application and native library for one target architecture using
the existing Docker setup. Verify installation and startup on the actual
gateway. Bound logs and event storage; if using a service, configure ordinary
automatic restart. Verify that the representative input rate fits the target's
memory and processing budget, including refits. This is a deployment check, not
a new performance-testing framework. Keep processing time distinct from the
time spent waiting for a ten-minute window to close.

\textbf{Where:} CLI, scoring, artefact loading, Docker/native build, and a short
installation/run guide. Share the existing scoring and adaptation logic.

\textbf{Done when:} one documented command processes successive windows on the
gateway, emits explainable decisions, performs guarded updates, restores state
after restart, and runs without continually growing memory or disk use. A small
integration check should confirm that the runtime matches the existing scorer
on identical windows and rejects malformed input.

\subsubsection*{Implementation breakdown for the gateway runtime}

Follow the items below in order. Each has a concrete change and a check that
shows when it is finished. They describe work to implement, not work already
verified on a gateway.

\textbf{Starting point on this branch.} Versioned detector state, atomic
activation, restart loading, a process lock, rollback, and bounded state history
already exist. Reuse them. Both conservative and legacy adaptation policies
also exist; retain access to the earlier methods for reproducible thesis
experiments. A legacy result must identify its policy. A run where every
candidate is rejected demonstrates rejection behaviour, not successful learning.
The edge work should preserve these behaviours while replacing the batch input
and output path.

\subsubsection*{4.1 Define exactly what the gateway receives and loads}

\textbf{Local implementation status, 15 October 2026:} implemented and checked.
Models now store a fingerprint-bound input contract; prepared state loads with a
partition-free edge configuration. The three bundled model/threshold pairs were
regenerated without changing alpha or numeric thresholds. Older saved-state
directories require fresh preparation. See \path{docs/edge-detector.md}.

\textbf{Change:} keep the current four counts. Document the target device/MAC,
category order, trained window duration, device filtering, and the rule that
SSDP packets are excluded from ordinary UDP counts. Store this input contract
with the trained model and validate it when loading the matching threshold and
saved detector state. Complete any missing compatible-input checks from the
first requirement before trusting a stream. An operator-supplied duration alone
cannot prove which duration trained an old model; regenerate artefacts if their
format changes.

Add a small edge configuration containing artefact/state paths and the expected
input and detector settings. It must not require lists of fitting, calibration,
or future capture dates. Keep those lists in saved artefacts as research
provenance. Separate shared inference-compatibility checks from research
partition checks in the existing loaders, including the saved-state loader;
retain device, input, backend, score-type, policy, and model--threshold checks.
Prepare the initial saved state on the development machine so that the gateway
does not need calibration captures to reconstruct its conservative count limit.

\textbf{Where:} \path{src/lm_idnet/config.py},
\path{src/lm_idnet/artifacts.py}, model/threshold/state schemas and fingerprints,
and a small example under \path{configs/}.

\textbf{Check:} a valid prepared detector loads without research data folders or
capture partitions. Wrong device, category order, duration, counting policy,
non-finite threshold, or mismatched pair fails before scoring. Existing research
commands retain their partition validation.

\subsubsection*{4.2 Share the existing one-window detector logic}

\textbf{Local implementation status, 15 October 2026:} implemented and checked.
The batch runner and edge Python API share \texttt{WindowDetector.process()}.
Twelve controlled mode/policy/score combinations match the recorded pre-refactor
behaviour, and all six historical legacy experiments reproduce their scores and
decisions. Gateway hardware validation remains outstanding.

\textbf{Change:} extract the per-window work from the existing batch loop into
one shared processing path. Let the batch runner continue to collect its reports,
while the gateway keeps only the active detector, bounded learning buffers,
and counters. Initialise the LM backend and refit estimator once. Replace uses
of the accumulated result-list length with a scored-window counter. Do not keep
an ever-growing event or update-attempt list in the gateway path.

Keep the same order: score with the current pair, decide whether the window can
be learned from, build and check an update when due, then save an accepted pair
before activating it. Preserve the original safety references and count limit
when restoring state. Support static, threshold-adaptive, and periodic-refit
modes with the existing conservative/legacy policy choices; conservative remains
the default. This refactor must not change their statistical rules.

\textbf{Where:} \path{src/lm_idnet/models/scoring_stage.py}, its existing callers
and focused scoring/adaptation tests. Use the current artefact helpers.

\textbf{Check:} before switching the batch runner, record a small sequence's
scores, decisions, update outcomes, alpha, and threshold. Both callers then
match those results for both policies, including an accepted update and a
rejected update. Learning buffers stay within their configured sizes.

\subsubsection*{4.3 Add the command that reads and writes one window at a time}

\textbf{Local implementation status, 15 October 2026:} implemented and checked.
The \texttt{run} command reads bounded UTF-8 JSON lines, validates chronology and
input semantics, and flushes each event immediately. Its logs rotate at 5 MiB
with two backups. Tests verify delivery before input closes, malformed-input
errors, adaptive updates, and state locking. Service-signal shutdown and gateway
hardware validation remain outstanding. See \path{docs/edge-detector.md}.

\textbf{Change:} add a \texttt{run} command to the existing CLI. The implemented
interface is:
\begin{verbatim}
lm-idnet run --config configs/edge_example.json \
  --state-dir artifacts/edge-state
\end{verbatim}
Use one JSON object per input line: \texttt{device\_id}, \texttt{source\_id}, and
a nested \texttt{window} using the existing \texttt{WindowRecord} fields
(\texttt{start\_utc}, \texttt{end\_utc}, \texttt{categories}, \texttt{counts},
\texttt{state}). The source identifier supplies the existing event's
\texttt{capture\_id} context; it need not contain a date or belong to a research
partition. Supply one valid example line in the run guide.

Validate each line before processing it: counts must be actual non-negative
integers within the numerical implementation's supported range, rather than
booleans, strings, rounded fractions, or overflowing totals. Check the device,
category order, exact duration, UTC times, and chronological non-overlap. Reject
duplicates and backward windows within a running stream. Gaps represent absent
coverage, not invented zero counts. Skip explicit missing windows without
scoring or learning; score genuine silent windows normally.

Choose and document a maximum line size, and enforce it while reading, before
an unbounded line can enter memory. Invalid JSON, non-finite numbers, incomplete
records, and over-limit lines stop with a line-numbered error on standard error
and a nonzero exit status. Standard output contains only the existing JSON
events, including scores and explanations. Flush each event as it is processed;
send status messages and update reasons to bounded logs. Read and write
synchronously so a slow consumer does not create a growing queue.

\textbf{Where:} \path{src/lm_idnet/cli.py}, a small runtime module if needed,
the shared input schemas, CLI tests, and
\path{docs/command-line-interface.md}.

\textbf{Check:} send one line while keeping input open and receive its complete
event before sending the next line. Malformed, oversized, wrong-device, and
wrong-duration lines fail clearly. Missing and silent records behave differently,
and every standard-output line parses as an event.

\subsubsection*{4.4 Connect restart, shutdown, and rollback to the live command}

\textbf{Change:} hold the existing state-directory lock for the runtime's
lifetime. Load the accepted pair at startup, and use the existing atomic save
for every promotion. Keep pre-commit failures from changing the active pair;
retain the existing stop-on-uncertain-commit behaviour. Exit cleanly on end of
input and normal service shutdown, releasing the lock and avoiding a half-finished
activation. Learning buffers and update counters can restart empty.

Document rollback as an operator action: stop the runtime, invoke the existing
rollback command with compatible configuration, then start again. Make that
command accept the edge configuration through the same compatibility checks.
Do not add a second state format or a hot-rollback mechanism.

\textbf{Important limit:} saved model state does not save the input position or
guarantee that every output was delivered exactly once. State clearly that a
replayed window can be scored and learned from again after restart. The input
producer must control replay; a durable queue or cursor is outside this minimum.

\textbf{Where:} the runtime/CLI lifecycle and the existing
\path{src/lm_idnet/artifacts.py} lock, save, load, and rollback helpers.

\textbf{Check:} accept an update, stop, and restart without the original research
captures; the accepted alpha and threshold return together. Verify rollback,
lock refusal for a second writer, shutdown, and a failed save using small
integration tests that reuse the existing persistence tests.

\subsubsection*{4.5 Add the missing false-positive rate to the existing evaluator}

\textbf{Change:} extend the existing labelled evaluation output with
\(\mathrm{FPR}=\mathrm{FP}/(\mathrm{FP}+\mathrm{TN})\). Use JSON
\texttt{null} when the denominator is zero. Explain in the output documentation
that recall is the true-positive rate; keep the existing accuracy, recall, and
confusion counts. Keep this as an offline utility. Neither labels nor metric
calculation belongs in the live scoring or learning loop.

\textbf{Where:} \path{src/lm_idnet/evaluation/evaluation_stage.py}, its result
schema if applicable, evaluation tests, and CLI output documentation.

\textbf{Check:} a tiny hand-counted example gives the expected FPR; no benign
labels gives \texttt{null}, and no false positives gives zero. Existing metrics
remain unchanged and the output contains no NaN or infinity.

\subsubsection*{4.6 Build one image that starts without compiling anything}

\textbf{Change:} name the intended Linux gateway, operating system, and CPU
architecture. Build the existing application and native LM library for that
architecture. Compile during image creation, install the normal package without
the test extra, and launch the installed command directly at startup. Avoid
introducing a new dependency or packaging framework.

Provide a deployment invocation that runs the already-built image, keeps
standard input open (\texttt{docker run -i}), and forwards events unchanged.
The current development runner rebuilds on every invocation and its entrypoint
prints compilation messages; the deployment path must start without a build,
dependency download, or status text on event output. Keep configuration and
initial artefacts read-only, mount state persistently with the runtime user's
write permission, and forward stop signals to the application. The development
runner can keep its build convenience.

\textbf{Where:} \path{Dockerfile}, \path{scripts/docker-entrypoint.sh}, and the
deployment invocation alongside \path{scripts/run.sh} or in the run guide.

\textbf{Check:} start the prepared image with networking disabled and process a
fixture through its standard input. The native library loads on the stated
architecture, output is clean JSON lines, and state survives container removal
and restart through the persistent mount.

\subsubsection*{4.7 Put explicit limits on unattended operation}

\textbf{Change:} supply one ordinary service setup with automatic restart on
failure. Connect it to the chosen local window producer and a bounded event
destination. State who restarts or reconnects that producer; the detector cannot
recreate windows that were never supplied. End of a finite input file is a clean
exit, so a replay fixture must not cause an endless service restart loop.

Replace the runtime's unbounded file logging with rotation or the selected
service's bounded journal. State concrete maximum sizes and retained file counts
for application logs, container logs, and event output. An unrestricted shell
append to an event file does not meet this requirement. Document the policy
when storage is full and when an output consumer stops reading. Reuse the
existing active/previous snapshots and last 100 update-history entries; no
additional archive of all versions is needed.

\textbf{Where:} CLI logging configuration, one service/deployment configuration,
and the installation/run guide.

\textbf{Check:} force log/event rotation with small test limits; storage remains
within the stated retention bounds. A process failure restarts into accepted
state, a requested stop stays stopped, and a slow output consumer does not cause
memory growth. A full-disk save failure preserves the valid pair as specified
by the existing commit handling.

\subsubsection*{4.8 Verify the complete path with a small repeatable replay}

\textbf{Change:} add a small prepared-window fixture and integration checks
through the real command. Cover normal traffic, silence, missing coverage, an
alert, an accepted update, and a rejected candidate. Compare the runtime with
the shared batch path on identical ordered windows and settings, for raw and
normalised scoring and the supported adaptation policies. Compare events,
model fingerprints, alpha, thresholds, and update reasons. Do not loosen
conservative limits merely to make a demonstration update happen; use a
controlled valid candidate for the correctness check.

Exercise bad JSON, a truncated record, oversized input, invalid counts,
non-finite numbers, wrong input semantics, duplicate/overlapping timestamps,
and valid gaps. Include the first-event-before-end-of-input check, graceful
shutdown, restart after promotion, and rollback. Repeat a bounded valid sequence
with increasing timestamps to confirm that result and attempt histories are not
accumulated. Reuse existing fixtures and persistence checks rather than building
a new testing framework.

\textbf{Where:} existing scoring, CLI, artefact, and evaluation tests plus a
small runtime integration test and input fixture.

\textbf{Check:} relevant existing tests and the added integration checks pass.
The deployed-image smoke check also passes with the same fixture. Controlled
update tests prove correctness; they do not establish that a real traffic run
will produce a useful conservative update.

\subsubsection*{4.9 Run it on the gateway and write the handover instructions}

\textbf{Change:} install the prepared image on the actual target. Replay a
representative workload that includes update attempts and refits. Record the
hardware, OS, CPU architecture, source revision, image/dependency versions,
detector configuration, input rate, and window duration. Measure startup time,
complete window-to-event processing latency, throughput, CPU use, peak process
memory, and state/log/event disk use. Include refit latency and watch memory and
storage over a long enough replay to reveal continuing growth. Reuse small
existing timing helpers where suitable.

Compare the measurements with the target's available memory/disk and required
input rate, including headroom for an update. Report ordinary and worst observed
processing times and any pauses from refitting. The wait for a ten-minute
traffic window to close is separate from processing latency. Development-host
measurements alone do not verify gateway deployment.

Write a short guide with the exact build/install/start/stop/restart/rollback
commands, initial state preparation, mounts and permissions, one working
configuration, one example input line, expected event fields, storage limits,
and the measured operating limits. Identify input as supplied windows unless a
packet collector was actually delivered. Preserve raw scores and name the
adaptation policy in any thesis demonstration.

\textbf{Where:} a short gateway run guide under \path{docs/}, the small replay
fixture, and a measurement report under \path{reports/}.

\textbf{Check:} another person can follow the guide on the named gateway, score
successive windows, inspect accepted/rejected updates, restart, and roll back.
Claim gateway validation only after this run; until then, report the remaining
hardware check explicitly. No web dashboard, HTTP endpoint, new traffic features,
or live packet collector is required for this supplied-window deliverable.

\clearpage
\section{The stopping point}

The minimum detector is complete when another person can:
\begin{enumerate}
  \item Train and calibrate a normal model through the existing commands.
  \item Install the runtime on the stated gateway and supply count windows.
  \item Receive scores, anomaly decisions, and the existing explanations.
  \item Observe a permitted adaptive update and a safely rejected update.
  \item Restart or roll back without losing a valid model--threshold pair.
\end{enumerate}
Hand over the source revision, one working configuration, initial model and
threshold, installation/run instructions, and the relevant passing correctness
checks. \textbf{Stop there.} These four tasks deliver the promised software.

\section*{Source basis}

\small
Original LM-IDNet brief, page 2:
\path{resouces/Proposte_di_Ricerca_basate_su_Dirichlet_Multinomial.pdf}.
Guarded adaptation and existing explanations:
\path{docs/thesis_proposal_and_system_architecture.tex}, stages 8--9.
Current command behaviour: \path{docs/command-line-interface.md}.
Implementation checked on the original reference commit:
\path{src/lm_idnet/models/schemas.py}, \path{src/lm_idnet/artifacts.py},
\path{src/lm_idnet/models/calibration_stage.py},
\path{src/lm_idnet/models/scoring_stage.py}, and
\path{src/lm_idnet/evaluation/evaluation_stage.py}.

Gateway breakdown checked against \texttt{feature/edge} on 15 October 2026:
the same artefact/scoring/schema files, \path{src/lm_idnet/config.py},
\path{src/lm_idnet/cli.py}, \path{src/lm_idnet/processing/schemas.py},
\path{src/lm_idnet/processing/categories.py}, \path{Dockerfile},
\path{scripts/docker-entrypoint.sh}, \path{scripts/run.sh}, and
\path{docs/edge-implementation-guideline.md}. The guideline's static pilot is
a useful intermediate check; this plan retains adaptive operation as the
completed detector's requirement. Its older adaptation-status and test-count
statements are not used as current status evidence.

\end{document}
